\begin{table}[ht]
\centering
\caption{Pricing-versus-sourcing reconciliation.}
\label{tab:utg_reconciliation}
\begin{threeparttable}
\small
\setlength{\tabcolsep}{4pt}
\begin{adjustbox}{max width=\textwidth}
\begin{tabular}{p{.34\linewidth}rrp{.22\linewidth}}
\toprule
Component & Log-points & Admin-vs-litigated (\%) & 95\% CI \\
\midrule
Observed UTG gap & \BPutgObsLogGap{} & \BPutgObsPct{} & [\BPutgObsCIlow{}, \BPutgObsCIhigh{}] \\
Mechanical quantity component & \BPutgMechLogGap{} & \BPutgMechanicalCone{} & [\BPutgMechCIlow{}, \BPutgMechCIhigh{}] \\
Within firm-buyer-item component & \BPutgWithinLogGap{} & \BPutgWithinFirmOffset{} & [\BPutgWithinCIlow{}, \BPutgWithinCIhigh{}] \\
Residual supplier-composition component & \BPutgCompLogGap{} & \BPutgCompositionResidual{} & [\BPutgCompCIlow{}, \BPutgCompCIhigh{}] \\
\bottomrule
\end{tabular}
\end{adjustbox}
\begin{tablenotes}[flushleft]\footnotesize
\item \textit{Notes:} Components are expressed in administrative-minus-litigated units. The decomposition uses \BPutgReconBootB{} PBU-cluster bootstrap draws for confidence intervals. The residual composition component is a reconciliation residual and should be read together with the direct winner-switching evidence in the main paper.
\end{tablenotes}
\end{threeparttable}
\end{table}
